\section*{Contoh}
\addcontentsline{toc}{section}{Contoh}

\begin{contoh}[Mengklasifikasikan lima tata bahasa]
Tentukan tipe setiap tata bahasa berikut.
\begin{enumerate}[label=(\alph*),leftmargin=*]
  \item \(S \to aS \mid bS \mid \varepsilon\).
  \item \(S \to aSb \mid \varepsilon\).
  \item \(S \to aSBC \mid aBC\), \(CB \to BC\), \(aB \to ab\),
        \(bB \to bb\), \(bC \to bc\), \(cC \to cc\).
  \item \(S \to AB\), \(AB \to BA\), \(A \to a\), \(B \to b\).
  \item \(S \to Sa \mid b\).
\end{enumerate}
\textbf{Jawaban.}
\begin{enumerate}[label=(\alph*),leftmargin=*]
  \item Tipe 3 (regular), karena semua aturan berbentuk \(A \to aB\)
        atau \(A \to \varepsilon\).
  \item Tipe 2 (bebas konteks), karena aturan berbentuk
        \(A \to \alpha\) dan menghasilkan \(\{a^n b^n\}\).
  \item Tipe 1 (sensitif konteks), karena aturan seperti
        \(CB \to BC\) bergantung pada konteks dan panjang string tidak
        berkurang. Tata bahasa ini membangkitkan
        \(\{a^n b^n c^n\}\).
  \item Tipe 0 (tanpa batas), karena aturan \(AB \to BA\) tidak
        memenuhi batasan sensitif konteks.
  \item Tipe 3 (regular), karena aturan berbentuk \(A \to Ba\) atau
        \(A \to a\), yaitu left-linear.
\end{enumerate}
\end{contoh}

\begin{contoh}[Tabel pemetaan Tipe 3--0]
Tabel berikut merangkum pemetaan yang harus diisi pada Sub-CPMK 1.4.
\begin{center}
\small
\begin{tabular}{@{}c l l@{}}
\toprule
\textbf{Tipe} & \textbf{Automata} & \textbf{Contoh bahasa} \\
\midrule
3 & DFA/NFA & \(a^{*}b^{*}\) \\
2 & PDA & \(\{a^n b^n \mid n \ge 0\}\) \\
1 & Otomata terbatas linear & \(\{a^n b^n c^n \mid n \ge 1\}\) \\
0 & Mesin Turing & bahasa yang dikenali mesin Turing \\
\bottomrule
\end{tabular}
\end{center}
Setiap baris harus konsisten: contoh bahasa pada kolom terakhir harus
benar-benar dikenali oleh automata pada kolom kedua.
\end{contoh}

\begin{contoh}[Dari ekspresi reguler ke otomata]
Ekspresi reguler \((a|b)^{*}abb\) menyatakan bahasa seluruh string
atas \(\{a,b\}\) yang berakhir dengan \(abb\). Otomata hingga untuk
bahasa ini memiliki empat state yang mengingat prefiks terpanjang
yang cocok: \(\varepsilon\), \(a\), \(ab\), dan \(abb\). State
terakhir adalah state penerima. Setiap simbol baru memperbarui state
sesuai prefiks terpanjang, sehingga otomata hanya memerlukan memori
konstan.
\end{contoh}
